% Generated draft: requires mathematical and source review.
\documentclass[11pt]{article}
\usepackage[margin=25mm]{geometry}
\usepackage{fontspec}
\setmainfont{Latin Modern Roman}
\usepackage{amsmath,amssymb,mathtools,unicode-math}
\setmathfont{Latin Modern Math}
\usepackage{enumitem,needspace,xurl,hyperref}
\setlength{\parindent}{0pt}
\setlength{\parskip}{5pt}
\setlength{\emergencystretch}{3em}
\newcommand{\N}{\mathbb{N}}
\newcommand{\Z}{\mathbb{Z}}
\newcommand{\Q}{\mathbb{Q}}
\newcommand{\R}{\mathbb{R}}
\newcommand{\C}{\mathbb{C}}
\newcommand{\F}{\mathbb{F}}
\newcommand{\A}{\mathbb{A}}
\newcommand{\PP}{\mathbb{P}}
\newcommand{\E}{\mathbb{E}}
\newcommand{\eps}{\varepsilon}
\newcommand{\dd}{\,\mathrm{d}}
\begin{document}
\section*{UnsolvedMath ID: 30002748}
\textbf{Source problem code:} catalogue entry
\subsection{Definitions and mathematical statement}
\textbf{Source question.} The supplied record asks: ``What are the points of difficulty in the mathematical situation, and how is perceived difficulty treated?''

\textbf{Unresolved definitions and scope.} The record does not specify the mathematical situation, the participants whose perceptions are considered, or the educational or other setting. It supplies no definition of a point of difficulty, no criterion for perceived difficulty, and no meaning or procedure for how that difficulty is ``treated.'' In particular, perceived difficulty cannot silently be identified with an objective mathematical measure of difficulty.

The question is therefore retained as a research question about mathematical activity and its interpretation. The supplied evidence does not determine mathematical objects, domains, hypotheses, boundary conditions, quantifiers, or a criterion for a complete answer. A precise quantified mathematical proposition cannot be extracted without additional source context and definitions.

\textbf{Supplied status qualification.} The registry labels the problem open. This is a supplied curation label, not an independently verified mathematical fact.
\subsection{Short English statement}
The question asks where difficulty arises in a mathematical situation and how perceived difficulty is addressed. The situation, participants, meaning of difficulty, and meaning of treatment are unspecified, so a precise mathematical statement requires further source context.
\subsection{Sources}
\begin{itemize}
\item Supplied local UnsolvedMath registry record, permanent ID \(30002748\), title: \emph{Perceived Difficulty in Mathematical Activity}; exact source problem code: \texttt{catalogue entry}. The record gives proposed year 2014 and names no individual proposer. Its statement status is supplied as ``exact,'' with the qualification that the wording was retained byte-for-byte because it was internally readable and complete.
\item Supplied report citation: \emph{Mathematics in Undergraduate Study Programs: Challenges for Research and for the Dialogue between Mathematics and Didactics of Mathematics} (2015), Oberwolfach Reports, pp.\ 3103--3175, DOI 10.4171/owr/2014/56. \url{https://doi.org/10.4171/owr/2014/56}. This is also the literature-source link labeled ``original OWR report.'' The supplied background describes Oberwolfach Reports as documenting mathematical workshops and their open-problem discussions.
\item Supplied literature source: ``Zazkis and Zazkis, proof-comprehension study.'' \url{https://doi.org/10.1007/s10763-014-9595-0}. No fuller bibliographic citation is supplied.
\item Literature assessment dated \textbf{2026-08-21}: supplied dated curation, not independently verified. It labels the problem open and states that proof-script studies identify selected points of difficulty and how participants address them, but do not establish a general theory answering how perceived mathematical difficulty is treated across settings. The background describes this as dated literature triage derived from the supplied OWR curation and calls for independent verification of time-sensitive status and cited claims.
\end{itemize}
\end{document}
